\clearpage
\begin{examplebox}
  \textbf{\textcolor{CExample}{例一：从奇函数直接读出常数项（基础）}}
\begin{center}\begin{tikzpicture}[font=\small,>=stealth,line join=round]
\draw[gray!40,rounded corners=3pt] (0,0) rectangle (14.2,8.5);\node at (7.1,7.9) {易错三：消去常数项，还要消去偶次项};
\node at (3.5,6.85){$y=x^3+x^2$：$r=0$，仍不够};\node at (10.6,6.85){$y=x^3$：还需让 $p=0$};\begin{scope}[shift={(3.5,3.3)},x=1cm,y=1cm]\draw[->,gray] (-2,0)--(2,0) node[right]{$x$};\draw[->,gray] (0,-1.6)--(0,3.0) node[above]{$y$};\draw[blue,thick,domain=-1.6:1.1,samples=150] plot(\x,{(\x)^3+(\x)^2});\fill[orange!85!black] (1,2) circle(2.2pt) node[above right=5pt] {$(1,2)$};\fill[orange!85!black] (-1,0) circle(2.2pt) node[below left=5pt] {$(-1,0)$};\fill[orange!85!black] (0,0) circle(2.2pt);\end{scope}\begin{scope}[shift={(10.5,3.3)},x=1cm,y=1cm]\draw[->,gray] (-2,0)--(2,0) node[right]{$x$};\draw[->,gray] (0,-1.6)--(0,3.0) node[above]{$y$};\draw[blue,thick,domain=-1.1:1.4,samples=150] plot(\x,{(\x)^3});\fill[orange!85!black] (1,1) circle(2.2pt) node[above right=5pt] {$(1,1)$};\fill[orange!85!black] (-1,-1) circle(2.2pt) node[below left=5pt] {$(-1,-1)$};\fill[orange!85!black] (0,0) circle(2.2pt);\end{scope}
\node[text width=12.8cm,align=center] at (7.1,.65) {图取 $q=0$，左图 $p=1,r=0$；右图 $p=r=0$。两图同尺度。};
\end{tikzpicture}\end{center}



  \textbf{题目：}已知 $f(x)=x^3+px^2+qx+r$ 是定义在全体实数上的奇函数，求 $r$，并判断 $p$ 需要满足什么条件。

  \textbf{最快做法：}因为定义域包含 $0$，而且 $f$ 是奇函数，
  \[
    f(0)=r=0,\qquad\boxed{r=0}.
  \]
  但\textbf{做到这里还不够}：$f(0)=0$ 只能确定常数项，不能保证整个函数是奇函数。
  把 $x$ 与 $-x$ 的函数值相加：
  \[
    f(x)+f(-x)=2px^2+2r.
  \]
  奇函数必须对所有 $x$ 满足这个和为零。已知 $r=0$，
  取 $x=1$ 得 $2p=0$，所以
  \[
    \boxed{p=0},\qquad q\text{ 可以是任意实数}.
  \]
  代回后 $f(x)=x^3+qx$，确实满足 $f(-x)=-f(x)$。

  \medskip
  
\end{examplebox}
\clearpage
\begin{examplebox}
\textbf{\textcolor{CExample}{例二：只需看常数项，不限制奇次项（对照）}}
\begin{center}\begin{tikzpicture}[font=\small,>=stealth,line join=round]
\draw[gray!40,rounded corners=3pt] (0,0) rectangle (14.2,8.5);\node at (7.1,7.9) {例二：奇次项可以组合，常数项必须为零};
\begin{scope}[shift={(4.2,3.8)},x=1.7cm,y=0.65cm]\draw[->,gray] (-1.8,0)--(1.8,0) node[right]{$x$};\draw[->,gray] (0,-4)--(0,4) node[above]{$y$};\draw[blue,thick,domain=-1.05:1.05,samples=150] plot(\x,{(\x)^5+(\x)^3+\x});\fill[orange!85!black] (1,3) circle(2.2pt) node[above right=5pt] {$(1,3)$};\fill[orange!85!black] (-1,-3) circle(2.2pt) node[below left=5pt] {$(-1,-3)$};\fill[orange!85!black] (0,0) circle(2.2pt) node[below right=5pt] {$O$};\end{scope}\node[align=left,anchor=west] at (8.1,4.2){图取 $f(x)=x^5+x^3+x$\\[11pt]每一项换 $-x$ 后都变号\\[11pt]任意实系数的奇次项之和\\[11pt]仍满足 $f(-x)=-f(x)$};
\node[text width=12.8cm,align=center] at (7.1,.65) {图取 $a=b=c=1,d=0$；一般情形看下方代数验证。};
\end{tikzpicture}\end{center}



  \textbf{题目：}已知 $f(x)=ax^5+bx^3+cx+d$ 是定义在全体实数上的奇函数，求 $d$，并说明 $a,b,c$ 的取值。

  \textbf{解：} $f(0)=d=0$，所以 $\boxed{d=0}$。
  此时
  \[
    f(-x)=-ax^5-bx^3-cx=-f(x).
  \]
  可见不论 $a,b,c$ 取什么实数，都能满足奇函数的规则。
  \textbf{答案：}$d=0$，$a,b,c$ 可为任意实数。

  \medskip
  
\end{examplebox}
\clearpage
\begin{examplebox}
\textbf{\textcolor{CExample}{例三：奇函数求参，再结合给定函数值（高考常见）}}
\begin{center}\begin{tikzpicture}[font=\small,>=stealth,line join=round]
\draw[gray!40,rounded corners=3pt] (0,0) rectangle (14.2,8.5);\node at (7.1,7.9) {例三：先确定原点，再用给定函数值求参数};
\begin{scope}[shift={(4.2,3.8)},x=1.7cm,y=0.85cm]\draw[->,gray] (-1.8,0)--(1.8,0) node[right]{$x$};\draw[->,gray] (0,-3.1)--(0,3.1) node[above]{$y$};\draw[blue,thick,domain=-1.1:1.1,samples=150] plot(\x,{(\x)^3+\x});\fill[orange!85!black] (1,2) circle(2.2pt) node[above right=5pt] {$(1,2)$};\fill[orange!85!black] (-1,-2) circle(2.2pt) node[below left=5pt] {$(-1,-2)$};\fill[orange!85!black] (0,0) circle(2.2pt) node[below right=5pt] {$O$};\end{scope}\node[align=left,anchor=west] at (8.3,4.1){求参后 $f(x)=x^3+x$\\[11pt]$a=c=0,\ b=1$\\[11pt]已知 $f(1)=2$\\[11pt]同时必有 $f(-1)=-2$};
\node[text width=12.8cm,align=center] at (7.1,.65) {图像画的是求参后的实际函数；正负两点关于原点对称。};
\end{tikzpicture}\end{center}



  \textbf{题目：}已知 $f(x)=x^3+ax^2+bx+c$ 为定义在全体实数上的奇函数，且 $f(1)=2$，求 $f(0)$ 和 $b$。

  \textbf{第一步：}奇函数且 $0$ 有定义，直接得
  \[
    f(0)=c=0.
  \]
  \textbf{第二步：}比较 $f(x)$ 与 $f(-x)$，
  或直接把两式相加：
  \[
    f(x)+f(-x)=2ax^2+2c=0.
  \]
  已知 $c=0$，取 $x=1$ 得 $a=0$。

  \textbf{第三步：}代入 $f(1)=2$：
  \[
    1+a+b+c=2
    \quad\Longrightarrow\quad 1+b=2
    \quad\Longrightarrow\quad b=1.
  \]
  \textbf{答案：}$\boxed{f(0)=0,\ b=1}$。

  \medskip
  
\end{examplebox}
\clearpage
\begin{examplebox}
\textbf{\textcolor{CExample}{例四：分式函数求参，先用 $f(0)=0$ 省步骤（进阶）}}
\begin{center}\begin{tikzpicture}[font=\small,>=stealth,line join=round]
\draw[gray!40,rounded corners=3pt] (0,0) rectangle (14.2,8.5);\node at (7.1,7.9) {例四：分母无零点，求参后整个图像关于原点对称};
\begin{scope}[shift={(4.2,3.8)},x=1.2cm,y=1.4cm]\draw[->,gray] (-2.5,0)--(2.5,0) node[right]{$x$};\draw[->,gray] (0,-1.9)--(0,1.9) node[above]{$y$};\draw[blue,thick,domain=-2.2:2.2,samples=150] plot(\x,{3*\x/(1+(\x)^2)});\fill[orange!85!black] (1,1.5) circle(2.2pt) node[above right=5pt] {$(1,\frac32)$};\fill[orange!85!black] (-1,-1.5) circle(2.2pt) node[below left=5pt] {$(-1,-\frac32)$};\fill[orange!85!black] (0,0) circle(2.2pt) node[below right=5pt] {$O$};\end{scope}\node[align=left,anchor=west] at (8.3,4.1){求参后 $f(x)=\dfrac{3x}{1+x^2}$\\[14pt]$a=2,\ b=1$\\[11pt]分母始终大于 $0$\\[11pt]原点存在\\[11pt]成对点高度互为相反数};
\node[text width=12.8cm,align=center] at (7.1,.65) {图中两点可直接代回验算；单有原点仍不能代替对所有 $x$ 的验证。};
\end{tikzpicture}\end{center}



  \textbf{题目：}已知
  \[
    f(x)=\frac{(a-2)x^2+3x+b-1}{x^2+1}
  \]
  是奇函数，求 $a,b$。

  \textbf{第一步：}因为分母 $x^2+1>0$，定义域是全体实数，包含 $0$。
  由 $f(0)=0$ 得
  \[
    \frac{b-1}{1}=0\quad\Longrightarrow\quad \boxed{b=1}.
  \]
  \textbf{第二步：}代入 $b=1$，
  \[
    f(x)=\frac{(a-2)x^2+3x}{x^2+1}.
  \]
  奇函数应满足 $f(x)+f(-x)=0$，而
  \[
    f(x)+f(-x)=\frac{2(a-2)x^2}{x^2+1}.
  \]
  为使每个实数 $x$ 都满足，取 $x=1$ 即得 $a-2=0$，
  所以 $\boxed{a=2}$。

  \textbf{检查：}此时 $f(x)=\dfrac{3x}{x^2+1}$，
  直接代入可见 $f(-x)=-f(x)$，确实为奇函数。
  \textbf{答案：}$\boxed{a=2,\ b=1}$。

  \medskip
  
\end{examplebox}
\clearpage
\begin{examplebox}
\textbf{\textcolor{CExample}{例五：奇函数却没有 $f(0)$（防止漏条件）}}
\begin{center}\begin{tikzpicture}[font=\small,>=stealth,line join=round]
\draw[gray!40,rounded corners=3pt] (0,0) rectangle (14.2,8.5);\node at (7.1,7.9) {易错一：关于原点对称，不等于图像一定包含原点};
\node at (3.5,6.85){$y=x^3$：$0$ 可以代入};\node at (10.6,6.85){$y=1/x$：$0$ 不能代入};\begin{scope}[shift={(3.5,3.7)},x=1cm,y=1cm]\draw[->,gray] (-2.2,0)--(2.2,0) node[right]{$x$};\draw[->,gray] (0,-2.2)--(0,2.2) node[above]{$y$};\draw[blue,thick,domain=-1.25:1.25,samples=150] plot(\x,{(\x)^3});\fill[orange!85!black] (0,0) circle(2.2pt) node[below right=5pt] {$O$};\end{scope}\begin{scope}[shift={(10.5,3.7)},x=1cm,y=1cm]\draw[->,gray] (-2.2,0)--(2.2,0) node[right]{$x$};\draw[->,gray] (0,-2.2)--(0,2.2) node[above]{$y$};\draw[blue,thick,domain=-2:-0.5,samples=150] plot(\x,{1/\x});\draw[blue,thick,domain=0.5:2,samples=150] plot(\x,{1/\x});\fill[orange!85!black] (1,1) circle(2.2pt) node[above right=5pt] {$(1,1)$};\fill[orange!85!black] (-1,-1) circle(2.2pt) node[below left=5pt] {$(-1,-1)$};\node[below right=5pt] at (0,0){$O$};\end{scope}\node at (3.5,1.25){是奇函数，且 $f(0)=0$};\node at (10.6,1.25){是奇函数，但 $f(0)$ 不存在};
\node[text width=12.8cm,align=center] at (7.1,.65) {两图横纵同尺度；右图原点仅是坐标轴交点，不是函数图像上的点。};
\end{tikzpicture}\end{center}



  \textbf{题目：}函数 $f(x)=\dfrac1x$ 是不是奇函数？能否得到 $f(0)=0$？

  \textbf{解：}它的定义域是 $x\ne0$，而且对每个允许的 $x$ 都有
  \[
    f(-x)=-\frac1x=-f(x).
  \]
  所以它\textbf{是奇函数}，但 $x=0$ 不能代入，
  \textbf{$f(0)$ 没有定义，绝不能写成 $0$}。

  \medskip
  
\end{examplebox}
\clearpage
\begin{examplebox}
\textbf{\textcolor{CExample}{例六：关于其他点中心对称，也能一秒求值（拓展）}}
\begin{center}\begin{tikzpicture}[font=\small,>=stealth,line join=round]
\draw[gray!40,rounded corners=3pt] (0,0) rectangle (14.2,8.5);\node at (7.1,7.9) {把原点换成 $(a,b)$：中心横坐标处只能取中心高度};
\begin{scope}[shift={(1.6,1.8)},x=1.4cm,y=0.8cm]\draw[->,gray] (-0.2,0)--(4.2,0) node[right]{$x$};\draw[->,gray] (0,-0.3)--(0,6.1) node[above]{$y$};\draw[blue,thick,domain=0.7:3.3,samples=150] plot(\x,{(\x-2)^3+3});\draw[gray,dashed] (2,0)--(2,5.7);\draw[gray,dotted] (0,3)--(3.9,3);\fill[orange!85!black] (2,3) circle(2.2pt) node[above left=5pt] {$S(2,3)$};\fill[orange!85!black] (1,2) circle(2.2pt) node[below right=5pt] {$(1,2)$};\fill[orange!85!black] (3,4) circle(2.2pt) node[above right=5pt] {$(3,4)$};\node[below=5pt] at (2,0){$2$};\node[left=5pt] at (0,3){$3$};\end{scope}\node[align=left,anchor=west] at (8.4,4.2){图取 $f(x)=(x-2)^3+3$\\[11pt]对称中心为 $(2,3)$\\[11pt]$f(2+t)+f(2-t)=6$\\[11pt]$x=2$ 有定义，所以 $f(2)=3$};
\node[text width=12.8cm,align=center] at (7.1,.65) {对称点 $(1,2)$、$(3,4)$ 的中点正是 $(2,3)$；一般情形仍要先确认 $x=a$ 有定义。};
\end{tikzpicture}\end{center}



  \textbf{题目：}函数 $f$ 在 $x=2$ 处有定义，且对每个允许的 $t$ 满足
  \[
    f(2+t)+f(2-t)=6.
  \]
  求 $f(2)$。

  \textbf{解：}取 $t=0$，得
  \[
    f(2)+f(2)=6
    \quad\Longrightarrow\quad\boxed{f(2)=3}.
  \]
  \textbf{理解：}这相当于图像关于点 $(2,3)$ 中心对称，
  对称中心所在的竖线上如果有图像点，就只能是这个中心本身。
  与本节“奇函数若在 $0$ 处有定义，则 $f(0)=0$”是同一个道理。
\end{examplebox}
